\section{Анализ предметной области разрабатываемого веб-приложения}

Современный учебный процесс в техническом вузе характеризуется интенсивной нагрузкой, большим числом дисциплин, лабораторных работ и проектов, а также сложной структурой расписания. Для РТУ МИРЭА характерны недельная цикличность (чётные и нечётные недели), частые изменения пар и распределение занятий по корпусам. Студенты вынуждены вести параллельно несколько источников данных: официальное расписание, личные заметки о дедлайнах, напоминания о консультациях и контрольных мероприятиях. При отсутствии единой системы это приводит к пропускам сроков и неэффективному планированию времени.

\subsection{Проблема управления дедлайнами и привязки к академическим событиям}

Большинство студентов ведут учёт сроков в заметках или календарях, не имеющих связи с расписанием. Такая организация данных не учитывает:
\begin{itemize}
    \item чётность недели и фактическое наличие пары (например, занятия проходят только раз в две недели);
    \item изменения расписания и переносы занятий;
    \item зависимость дедлайна от времени конкретной пары (например, сдача лабораторной работы до начала занятия по соответствующей дисциплине).
\end{itemize}
Отсутствие связи между расписанием и задачами вынуждает студента вручную переносить дедлайны, что увеличивает вероятность ошибок.

\subsection{Роль автоматизации и интеграции с расписанием}

Решением выступает автоматическая интеграция с данными расписания РТУ МИРЭА через выделенный сервис \texttt{rtu-mirea-schedule}. В результате приложение может:
\begin{itemize}
    \item получить список дисциплин для выбранной группы;
    \item показывать расписание на календарной сетке;
    \item автоматически предзаполнять дедлайн в момент создания задачи (по времени пары);
    \item визуально связывать задачи с конкретными предметами.
\end{itemize}
Тем самым сокращается число ручных действий и повышается точность планирования.

\subsection{Сравнительный анализ аналогов}

Для обоснования целесообразности разработки собственной системы выполнен обзор существующих решений. Критерии сравнения: наличие привязки к расписанию вуза, ориентация на учебные задачи и возможность интеграции с российскими сервисами.

\begin{table}[h!]
\centering
\caption{Сравнительный анализ систем управления задачами}
\begin{tabular}{|p{4.2cm}|p{3.2cm}|p{3.2cm}|p{3.2cm}|}
\hline
\textbf{Критерий} & \textbf{Todoist} & \textbf{Trello} & \textbf{MyStudyLife} \\ \hline
Интеграция с РТУ МИРЭА & Отсутствует & Отсутствует & Отсутствует \\ \hline
Фокус на учебных задачах & Низкий & Средний & Высокий \\ \hline
Привязка к парам & Нет & Нет & Частично (ручной ввод) \\ \hline
Поддержка расписания по группам & Нет & Нет & Нет \\ \hline
\end{tabular}
\end{table}

\begin{enumerate}
    \item \textbf{Todoist} --- универсальный таск-менеджер, перегруженный бизнес-функциями, без поддержки расписаний и учебных циклов.
    \item \textbf{Trello} --- инструмент Kanban, удобен для групповых проектов, но неудобен для ежедневного контроля дедлайнов по расписанию.
    \item \textbf{MyStudyLife} --- ориентирован на студентов, но требует полностью ручного ввода расписания и не учитывает специфику российских вузов.
\end{enumerate}

Таким образом, разработка специализированного веб-приложения, объединяющего расписание РТУ МИРЭА и персональные дедлайны студента, является актуальной и практически востребованной задачей.
